% !PRINTED 655
\weberpage{655}
%<<CONT
associirt. Ist also $\sigma^k$ die in $\omega$ enthaltene Potenz von $\sigma$, so setzen wir
\[
\omega = \sigma^k \omega',
\]
worin $\omega'$ mit der Primzahl $q$ theilerfremd ist. Die Bildung der Norm ergiebt nach 2.
\[
a^m = q^k b,
\]
worin $b$ eine rationale Zahl ist, die in einfachster Gestalt die Primzahl $q$ weder im Zähler noch im Nenner enthält. Daraus ergiebt sich aber, dass $k$ durch $m$ theilbar sein muss, und daraus das erste Resultat:

3. \gesperrt{Der Exponent der in $\omega$ enthaltenen Primzahl $\sigma$ ist immer durch $m$ theilbar.}

Es sei nun ferner $p$ eine von $q$ verschiedene Primzahl und $\mathfrak{p}_\xi$ ein Primfactor von $p$ im Körper $\Omega_m$. Wir lassen $\xi$ die im \S.\,172 in den verschiedenen Fällen näher bestimmte Zahlenreihe $\xi_1, \xi_2, \ldots, \xi_e$ durchlaufen, so dass
\weq{3}{p = \mathfrak{p}_{\xi_1}\mathfrak{p}_{\xi_2}\ldots\mathfrak{p}_{\xi_e}}
wird, und nehmen an, es sei $\mathfrak{p}_\xi^{k_\xi}$ die in $\omega$ enthaltene Potenz von $\mathfrak{p}_\xi$. Die in $\omega_n$ enthaltene Potenz von $\mathfrak{p}_{n\xi}$ ist dann $\mathfrak{p}_{n\xi}^{k_\xi}$, und wenn wir $n'$ aus der Congruenz
\weq{4}{n n' \equiv 1 \ (\mathrm{mod}\ m)}
bestimmen, so ist $\mathfrak{p}_\xi^{k_{n'\xi}}$ die in $\omega_n$ enthaltene Potenz von $\mathfrak{p}_\xi$. Demnach ist in
\weq{5}{\omega_n^{-1} \omega_1^n \text{ die Potenz } \mathfrak{p}_\xi^{n k_\xi - k_{n'\xi}}}
enthalten, und wegen 1. muss der Exponent dieser Potenz durch $m$ theilbar sein.

Wir erhalten also für jedes durch $q$ nicht theilbare $n$ die Congruenz
\weq{6}{n k_\xi \equiv k_{n'\xi} \ (\mathrm{mod}\ m),}
in der wir nun auch $\xi$ durch $n\xi$ ersetzen können, wodurch sich
\weq{7}{n k_{n\xi} \equiv k_\xi \ (\mathrm{mod}\ m)}
ergiebt. Nehmen wir hierin $\xi = 1$ und setzen dann $\xi, \xi'$ für $n, n'$ und $k$ für $k_1$, so folgt aus (7)
\weq{8}{k_\xi \equiv \xi' k \ (\mathrm{mod}\ m),}
worin nun $k$ für alle $\xi$ denselben Werth hat und $\xi, \xi'$ durch die Congruenz
\[
\xi \xi' \equiv 1 \ (\mathrm{mod}\ m)
\]
zusammenhängen.
